\section{通过生成元和关系定义的代数}

记号约定:$A$是交换环,$F,F^{\prime}$是$A$-代数,$J,H$是集合.

\begin{proposition}{赋值映射的性质}{}
	设$\left(x_{i}\right)_{i\in I}$是$F$的元素族,$f\in\Hom_{A-\mathsf{Alg}}\left(\Lib_{A}\left(I\right),F\right)$且$\forall i\in I\left(f\left(X_{i}\right)\right)=x_{i}$,则$f$是满射$\iff$ $\left(x_{i}\right)_{i\in I}$是$F$的生成系.
\end{proposition}

\begin{definition}{赋值映射}{}
	设$\mathbf{x}\coloneq\left(x_{i}\right)_{i\in I}$是$F$的元素族.
	\begin{enumerate}
		\item 我们称$\ev_{\mathbf{x}}:\Lib_{A}\left(I\right)\to F,X_{i}\mapsto x_{i}$是相对于$\mathbf{x}$的赋值映射.
		\item 对任一$U\in\Lib_{A}\left(I\right)$,称$U\left(\left(x_{i}\right)_{i\in I}\right)\coloneq\ev_{\mathbf{x}}\left(U\right)$是把不定元$X_{i}$替换为$x_{i}$得到的元素.
	\end{enumerate}
\end{definition}

\begin{proposition}{赋值映射和同态的兼容性}{}
	设$\left(x_{i}\right)_{i\in I}$是$F$的元素族,$g\in\Hom_{A-\mathsf{Alg}}\left(F,F^{\prime}\right)$,则对任一$U\in\Lib_{A}\left(I\right)$有$g\left(U\left(\left(x_{i}\right)_{i\in I}\right)\right)=U\left(\left(g\left(x_{i}\right)\right)_{i\in I}\right)$.
\end{proposition}

\begin{proposition}{复合赋值恒等式}{}
	对$\Lib_{A}\left(J\right)$的任一元素族$\left(H_{i}\right)_{i\in I}$和$F^{\prime}$的任一元素族$\left(y_{j}\right)_{j\in J}$,有$U\left(\left(H_{i}\right)_{i\in I}\right)\left(\left(y_{j}\right)_{j\in J}\right)=U\left(\left(H_{i}\left(\left(y_{j}\right)_{j\in J}\right)\right)_{i\in I}\right)$.
\end{proposition}

\begin{definition}{关系元和关系理想}{}
	设$\mathbf{x}\coloneq\left(x_{i}\right)_{i\in I}$是$F$的元素族.
	\begin{enumerate}
		\item 若$U\in\Lib_{A}\left(I\right)$且$\ev_{\mathbf{x}}\left(U\right)=0$,则称$U$是$\mathbf{x}$在$F$中的一个关系元(或关系).
		\item 称$\ker\left(\ev_{\mathbf{x}}\right)$是$\mathbf{x}$的关系理想.
	\end{enumerate}
\end{definition}

\begin{definition}{代数的展示}{}
	设$\mathbf{x}\coloneq\left(x_{i}\right)_{i\in I}$是$F$的生成系,$\Lib_{A}\left(J\right)$的元素族$\left(R_{j}\right)_{j\in J}$生成的双边理想是$\ker\left(\ev_{\mathbf{x}}\right)$.
	\begin{enumerate}
		\item $\left(\left(x_{i}\right)_{i\in I},\left(R_{j}\right)_{j\in J}\right)$称为$F$的一个展示.
		\item 对每个$i\in I$,称$x_{i}$是该展示的生成元.
		\item 对每个$j\in J$,称$R_{j}$是该展示的定义关系(或关系元).
	\end{enumerate}
\end{definition}

\begin{definition}{生成系和关系定义的泛代数}{}
	设$\mathfrak{r}$是$\Lib_{A}\left(I\right)$的元素族$\left(R_{j}\right)_{j\in J}$生成的双边代数理想.商代数$\Lib_{A}\left(I\right)/\mathfrak{r}$称为生成系$I$和关系族$\left(R_{j}\right)_{j\in J}$定义的泛代数.
\end{definition}

\begin{proposition}{泛代数的性质和泛性质}{}
	设$\mathfrak{r}$是$\Lib_{A}\left(I\right)$的元素族$\left(R_{j}\right)_{j\in J}$生成的双边代数理想.对诸$i\in I$,记$\bar{X}_{i}$是$X_{i}$在典范映射$\Lib_{A}\left(I\right)\twoheadrightarrow\Lib_{A}\left(I\right)/\mathfrak{r}$下的像.
	\begin{enumerate}
		\item $\left(\left(\bar{X}_{i}\right)_{i\in I},\left(R_{j}\right)_{j\in J}\right)$是$\Lib_{A}\left(I\right)/\mathfrak{r}$的展示.
		\item 设$F$的元素族$\left(x_{i}\right)_{i\in I}$满足$\forall j\in J\left(R_{j}\left(\left(x_{i}\right)_{i\in I}\right)=0\right)$,则存在唯一的$g\in\Hom_{A-\mathsf{Alg}}\left(\Lib_{A}\left(I\right)/\mathfrak{r},F\right)$使下图交换.
		\begin{center}
			\begin{tikzcd}
				{\Lib_{A}\left(I\right)} && {\Lib_{A}\left(I\right)/\mathfrak{r}} \\
				& F
				\arrow[two heads, from=1-1, to=1-3]
				\arrow[from=1-1, to=2-2]
				\arrow["{\exists!g}", dashed, from=1-3, to=2-2]
			\end{tikzcd}
		\end{center}
		\item 在(2)的条件下,$\left(\left(x_{i}\right)_{i\in I},\left(R_{j}\right)_{j\in J}\right)$是$F$的展示$\iff$ $g$是双射.
	\end{enumerate}
\end{proposition}

\begin{definition}{泛关系元/恒等式}{}
	设$S\in\Lib_{A}\left(H\right)$.若对$F$的每个元素族$\left(x_{h}\right)_{h\in H}$有$S\left(\left(x_{h}\right)_{h\in H}\right)=0$,则称$S$是$F$的一个泛关系元(或恒等式).
\end{definition}

\begin{example}{经典代数类的恒等式刻画}{}
	以$\left(X_{1}X_{2}\right)X_{3}-X_{1}\left(X_{2}X_{3}\right)$为泛关系元的代数是$A$-结合代数,以$X_{1}X_{2}-X_{2}X_{1}$为泛关系元的代数是$A$-交换代数.
\end{example}

\begin{definition}{恒等式定义的泛代数}{}
	设$\left(S_{k}\right)_{k\in K}$是$\Lib_{A}\left(H\right)$的元素族.记
	\begin{align*}
		T=\left\{S_{k}\left(\left(U_{h}\right)_{h\in H}\right)\in\Lib_{A}\left(I\right)\middle|k\in K,\left(U_{h}\right)_{h\in H}\text{是}\Lib_{A}\left(I\right)\text{的元素族}\right\},
	\end{align*}
	以$\langle T\rangle$表示$T$生成的双边理想.称$\Lib_{A}\left(I\right)/\langle T\rangle$为生成系$I$和泛关系元族$\left(S_{k}\right)_{k\in K}$定义的泛代数.
\end{definition}

\begin{proposition}{}{}
	记号同上.对诸$i\in I$,记$\bar{X}_{i}$是$X_{i}$在典范映射$\Lib_{A}\left(I\right)\twoheadrightarrow\Lib_{A}\left(I\right)/\langle T\rangle$下的像.
	\begin{enumerate}
		\item 对每个$k\in K$,$S_{k}$是$F$的泛关系元.
		\item 设$\left(x_{i}\right)_{i\in I}$是$F^{\prime}$的生成系,对每个$k\in K$,$S_{k}$是$F^{\prime}$的泛关系元,则存在唯一的$h\in\Hom_{A-\mathsf{Alg}}\left(F,F^{\prime}\right)$使得
		\begin{align*}
			\forall i\in I\left(h\left(\bar{X}_{i}\right)=x_{i}\right).
		\end{align*}
	\end{enumerate}
\end{proposition}














































